\documentclass{oxmathproblems}

\printanswers

\course{Understanding Analysis, 2nd Edition}
\sheettitle{Chapter 5 - Section 5.2 - Derivatives and the Intermediate Value Property} 

\begin{document}
\begin{questions}

%------------------------- Problem 1 -------------------------
\miquestion Supply proofs for part (i) and (ii) of Theorem 5.2.4.
\begin{solution}
  Part (i): $(f+g)'(c) = f'(c) + g'(c)$

  Proof:
  \begin{align*}
     (f+g)'(c) = \lim_{x \to c}\frac{f(x)+g(x)-f(c)-g(c)}{x-c} &= \lim_{x \to c}\left(\frac{f(x)-f(c)}{x-c} + \frac{g(x)-g(c)}{x-c}\right)\\ 
     &= f'(c)+g'(c)
  \end{align*}
  where the last step is justified because we know both limits exist, so we can apply the Algebraic Limit Theorem (Theorem 4.2.4).

  Part (ii): $(kf)'(c) = kf'(c)$ for all $k \in \textbf{R}$.
  
  Proof:
  \[(kf)'(c) = \lim_{x \to c} k\left( \frac{f(x)-f(c)}{x-c} \right) = k \lim_{x \to c}\frac{f(x)-f(c)}{x-c} = kf'(c)\]
\end{solution}

%------------------------- Problem 2 -------------------------
\miquestion Exactly one of the following requests is impossible. Decide which it is, and provide examples for the other three. In each case, assume the functions are defined
on all of \textbf{R}.
\begin{enumerate}[label=(\alph*)]
  \item Functions $f$ and $g$ not differentiable at zero but where $fg$ is differentiable at zero.
  \item A function $f$ not differentiable at zero and a function $g$ differentiable at zero where $fg$ is differentiable at zero.
  \item A function $f$ not differentiable at zero and a function $g$ differentiable at zero where $f+g$ is differentiable at zero.
  \item A function $f$ differentiable at zero but not differentiable at any other point.
\end{enumerate}

\begin{solution}
  \begin{enumerate}[label=(\alph*)]
    \item $f(x) = 0$ for $x \geq 0$ and $f(x)=1$ otherwise. $g(x)=1$ for $x \geq 0$ and $g(x) = 0$ otherwise. Both functions are discontinuous at zero hence not differentiable there, but
    $fg$ is the zero function.
    \item $f(x) = 0$ for $x \geq 0$ and $f(x)=1$ otherwise. $g(x)=0$ everywhere.
    \item Impossible. If $f+g$ were differentiable at zero, then $(f+g)-g = f$ must be differentiable at zero as well by Theorem 5.2.4.
    \item $f(x) = x^2$ for $x \in \textbf{Q}$ and $f(x) = 0$ otherwise.
    
    It's differentiable at zero because for $c= 0$, $\frac{f(x)-f(c)}{x-c} = x$ and $\lim_{x \to 0} x = 0$.

    It's not differentiable at $c \neq 0$. Consider a sequence $(x_{n}) \to c$. We can split it into two subsequences $(a_{n}) \in \textbf{Q}$ and $(b_{n}) \notin \textbf{Q}$. We have
    $(a_{n}) \to c^2$ and $(b_{n}) \to 0$. Since $c^2 \neq 0$ for $c \neq 0$, the two subsequences converge to different limits, hence $(x_{n})$ doesn't converge.
  \end{enumerate}
\end{solution}

%------------------------- Problem 3 -------------------------
\miquestion
\begin{enumerate}[label=(\alph*)]
  \item Use Definition 5.2.1 to produce the proper formula for the derivation of $f(x)=1/x$.
  \item Combine the result in part (a) with the Chain Rule to supply a proof for part (iv) of Theorem 5.2.4.
  \item Supply a direct proof of Theorem 5.2.4 (iv) by algebraically manipulating the difference quotient for $(f/g)$ in a style similar to the proof of Theorem 5.2.4 (iii).
\end{enumerate}

\begin{solution}
  \begin{enumerate}[label=(\alph*)]
    \item \[\frac{f(x)-f(c)}{x-c}=\frac{\frac{1}{x}-\frac{1}{c}}{x-c}=\frac{-(x-c)}{xc}\frac{1}{x-c}=-\frac{1}{xc}\]

    So we have $f'(c)=\lim_{x \to c}(-\frac{1}{xc})=-\frac{1}{c^2}$ for $c \neq 0$.

    \item We first derive a formula for the derivative of $h(x)=1/g(x)$.
    
    \[\frac{h(x)-h(c)}{x-c}=\frac{-(g(x)-g(c))}{g(x)g(c)}\frac{1}{x-c}=-\frac{1}{g(x)g(c)} \cdot \frac{g(x)-g(c)}{x-c}\]

    Taking the limit to $c$ and noting that the second term of the product is the definition of $g'(c)$
    \[h'(c) = \lim_{x \to c} -\frac{1}{g(x)g(c)} \cdot \frac{g(x)-g(c)}{x-c} = -\frac{g'(c)}{g(c)^2}\]

    Now apply the product rule to $(fh)'(c)$ to get the formula for $(f/g)'(c)$.

    ----- Simpler solution by chatgpt -----\\
    Derive the formula of $h'(x)$ by defining $h(x)=1/x$ then apply the chain rule to $h(g(x))$.

    \item
    \begin{align*}
       \frac{\frac{f}{g}(x) - \frac{f}{g}(c)}{x-c} &=  \frac{\frac{f}{g}(x) - \frac{f(x)}{g(c)} + \frac{f(x)}{g(c)} - \frac{f}{g}(c)}{x-c}\\
       &= \frac{f(x)}{x-c}\left(\frac{g(c)-g(x)}{g(x)g(c)}\right) + \frac{1}{g(c)(x-c)}\left(f(x)-f(c)\right)\\
       &= -\frac{f(x)}{g(x)g(c)}\left( \frac{g(x)-g(c)}{x-c} \right) + \frac{1}{g(c)} \left( \frac{f(x)-f(c)}{x-c} \right)
    \end{align*}

    Take the limit to $c$ to get
    \begin{align*}
      -\frac{f(c)g'(c)}{g(c)^2} + \frac{f'(c)}{g(c)} = \frac{f'(c)g(c)-f(c)g'(c)}{g(c)^2}
    \end{align*} 
  \end{enumerate}
\end{solution}

%------------------------- Problem 4 -------------------------
\miquestion Follow these steps to provide a slightly modified proof of the Chain Rule.
\begin{enumerate}[label=(\alph*)]
  \item Show that a function $h: A \to \textbf{R}$ is differentiable at $a \in A$ if and only if there exists a function $l : A \to \textbf{R}$ which is continuous at $a$ and satisfies
  \[h(x)-h(a) = l(x)(x-a) \qquad \text{for all } x \in A.\]
  \item Use this criterion for differentiability (in both directions) to prove Theorem 5.2.5.
\end{enumerate}

\begin{solution}
  \begin{enumerate}[label=(\alph*)]
    \item First we prove that if such a function $l$ exists then $h$ is differentiable at $a$. We have
    \[h'(a) = \lim_{x \to a}\frac{h(x)-h(a)}{x-a} = \lim_{x \to a}l(x) = l(a)\]

    Next, we prove that if $h$ is differentiable at $a$ then such a function $l$ exists. Define
    \[
      l(x) = 
      \begin{cases}
        \frac{h(x)-h(a)}{x-a} & x \neq a\\
        h'(a) & (x=a)
      \end{cases}
    \]

    $l$ is continuous at $a$ since $\lim_{x \to a}l(x) = \lim_{x \to a}\frac{h(x)-h(a)}{x-a} = h'(a) = l(a)$.

    $h(x)-h(a)=l(x)(x-a)$ for all $x \in A$ by definition.

    \item TODO.
  \end{enumerate}
\end{solution}

%------------------------- Problem 5 -------------------------
\miquestion Let $ l(x) = \begin{cases} x^a & x > 0\\0 & x \leq 0 \end{cases}$.
\begin{enumerate}[label=(\alph*)]
  \item For which values of $a$ is $f$ continuous at zero?
  \item For which values of $a$ is $f$ differentiable at zero? In this case, is the derivative function continuous?
  \item For which values of $a$ is $f$ twice-differentiable?
\end{enumerate}

\begin{solution}
  \begin{enumerate}[label=(\alph*)]
    \item $a > 0$. The limit at 0 from the left is 0. If $a < 0$ then the right limit diverges. If $a=0$ then the right limit is 1. But if $a > 0$ then the right limit is 0.
    \item From the direction of positive $x$, we have $\lim_{x \to 0^{+}}\frac{x^a-0^a}{x-0}=\lim_{x \to 0^{+}}x^{a-1} = 0^{a-1}$. This is only defined for $a > 1$, in which case it's 0 and
    agrees with the left limit, so the derivative is defined.

    The derivative is continuous, but seems to require us to use the fact that $l'(x)=ax^{a-1}$ for \textit{non-integer} $a$, which we have not proved yet.
    \item The next right-side derivative at zero is $\lim_{x \to 0^{+}} x^{a-2} = 0^{a-2}$, so we need $a > 2$.
  \end{enumerate}
\end{solution}

%------------------------- Problem 6 -------------------------
\miquestion Let $g$ be defined on an interval $A$ and let $c \in A$.
\begin{enumerate}[label=(\alph*)]
  \item Explain why $g'(c)$ in Definition 5.2.1 could have been given by
  \[g'(c) = \lim_{h \to 0}\frac{g(c+h) - g(c)}{h}\]
  \item Assume $A$ is open. If $g$ is differentiable at $c \in A$, show
  \[g'(c) = \lim_{h \to 0}\frac{g(c+h) - g(c-h)}{2h}\]
\end{enumerate}

\begin{solution}
  \begin{enumerate}[label=(\alph*)]
    \item Let $h = x-c$. Then we have
    \[g'(c) = \lim_{(x-c) \to 0} \frac{g(x)-g(c)}{x-c}\]
    And use the fact that $\lim_{(x-c) \to 0} = \lim_{x \to c}$.
    \item
    \begin{align*}
      \lim_{h \to 0}\frac{g(c+h)-g(c-h)}{2h} &= \lim_{h \to 0}\frac{g(c+h)-g(c)+g(c)-g(c-h)}{2h}\\
      &= \frac{1}{2}\lim_{h \to 0}\left( \frac{g(c+h)-g(c)}{h} + \frac{g(c)-g(c-h)}{h}\right)\\
      &= \frac{1}{2}(g'(c)+g'(c)) = g'(c)
    \end{align*}

    We used the fact that $g'(c) = \lim_{h \to 0}\frac{g(c)-g(c-h)}{h}$, which we'll now prove. Let $h=c-x$. Then we have
    \[\lim_{(c-x) \to 0}\frac{g(c)-g(x)}{c-x} = \lim_{(c-x) \to 0}\frac{g(x)-g(c)}{x-c}\]
    And use the fact that $\lim_{(c-x) \to 0} = \lim_{(x-c) \to 0} = \lim_{x \to c}$ (this is easy to prove with the usual $\epsilon-\delta$).

    One more thing to note is why the problem assumes $A$ is open. This is so that both $g(c-h)$ and $g(c+h)$ are defined.
    
    \textbf{--- Slight improvement by chatgpt ---}\\
    To prove $g'(c) = \lim_{h \to 0}\frac{g(c)-g(c-h)}{h}$, set $k=-h$. Then as $h \to 0$, also $k \to 0$, and
    \[\frac{g(c)-g(c-h)}{h}=\frac{g(c+k)-g(c)}{k}\]
    Then reuse part (a).
  \end{enumerate}
\end{solution}

%------------------------- Problem 7 -------------------------
\miquestion TODO
%------------------------- Problem 8 -------------------------
\miquestion TODO

%------------------------- Problem 9 -------------------------
\miquestion True / false with proof / counterexample.
\begin{enumerate}[label=(\alph*)]
  \item If $f'$ exists on an interval and is not constant, then $f'$ must take on some irrational values.
  \item If $f'$ exists on an open interval and there is some point $c$ where $f'(c) > 0$, then there exists a $\delta-$neighborhood $V_{\delta}(c)$ around $c$ in which
  $f'(x) > 0$ for all $x \in V_{\delta}(c)$.
  \item If $f$ is differentiable on an interval containing zero and if $\lim_{x \to 0}f'(x) = L$, then it must be that $L = f'(0)$.
\end{enumerate}

\begin{solution}
  \begin{enumerate}[label=(\alph*)]
    \item True, by Darboux Theorem and the fact that the irrationals are dense in $\textbf{R}$.
    \item 
    \textbf{--- Incorrect ---}\\
    True. Otherwise $f'(c)$ is an isolated point and that violates Darboux Theorem.

    \textbf{--- Corrected ---} \\
    The above is wrong because Darboux Theorem only guarantees Intermediate Value Property and not continuity (see Section 4.5).

    I would not have been able to come up with the counterexample myself. It's
    \[
      f(x) = 
      \begin{cases}
        x/2 + x^2\sin(1/x) & x \neq 0\\
        0 & x = 0
      \end{cases}
    \]
    The $x/2$ term makes $f'(0) = \frac{1}{2}$ (compare this with $f(x)=x^2+\sin(1/x)$ where $f'(0)$ would be 0).
    
    For $x \neq 0$, $f'(x)=\frac{1}{2} - \cos(1/x) + 2x\sin(1/x)$.
    For any $V_{\delta}(0)$ there's always an $x$ such that $\cos(1/x)=1$, which would make $f'(x)$ negative.
  
    \item 
    \textbf{--- Incorrect ---}\\
    We need the following lemma: if $\lim_{a \to b}f(a)$ is defined for some $b$ near $c$, and $\lim_{a \to c}f(a)$ is also defined,
    then $\lim_{b \to c}(\lim_{a \to b}f(a))=\lim_{a \to c}f(a)$.
  
    The proof is straightforward with $\epsilon-\delta$. There's a $\delta$ such that $|a-c| < \delta \implies |f(a)-L| < \epsilon$.
    Make $|b-c|< \delta/2$ and $|a-b| < \delta/2$.

    It's very important to note that $\lim_{a \to b}f(a)$ *must* be defined for $b$ near $c$ for this lemma to work. Otherwise,
    one counterexample is the Dirichlet function. $\lim_{a \to 0}f(a)=0$. But the limit is undefined for all $b \neq 0$ so the expression
    $\lim_{b \to 0}(\lim_{a \to b}f(a))$ is undefined.

    Now back to the original problem. $f$ is differentiable on an interval containing zero, so $f'(0)$ exists and also
    \[\lim_{x \to c}\frac{f(x)-f(c)}{x-c}\]
    exists for $c$ near $0$.

    We also have
    \[\lim_{c \to 0}\left( \lim_{x \to c}\frac{f(x)-f(c)}{x-c}\right)=L\]
    Now apply the lemma to conclude that
    \[\lim_{c \to 0}\left( \lim_{x \to c}\frac{f(x)-f(c)}{x-c}\right)=\lim_{x \to 0}\frac{f(x)-f(0)}{x}=f'(0)=L\]
  
    \textbf{--- Corrected ---}\\
    The nested limit lemma above only applies to functions with a single parameter. For a counterexample, consider
    $\lim_{c \to 0}(\lim_{x \to c}\frac{x}{c}) = \lim_{c \to 0}1 = 1$.
    But $\lim_{x \to 0}\frac{x}{c}$ is undefined.

    Instead we can use Darboux Theorem. Suppose for contradiction that $f'(0) \neq L$ and let $d = |f'(0)-L|$. By convergence to $L$, there exists $\delta$ such that
    $x \in V_{\delta}(0) \implies |f'(x)-L| < d/2$.
    
    But by Darboux, for such $x \in V_{\delta}(0)$ then the derivative must take every value between
    $f'(x)$ and $f'(0)$. In particular, there exists $y \in V_{\delta}(0)$ such that $|f'(y)-f'(0)| < d/2$, contradicting convergence to $L$ ($f'(y)$
    cannot be within $d/2$ of both $L$ and $f'(0)$).
  \end{enumerate}
\end{solution}

\end{questions}
\end{document}
